% State-specific — Circle Graphs
\section{Circle Graphs}

\begin{learningGoals}
\begin{itemize}
    \item Read and interpret circle (pie) graphs
    \item Relate sectors to fractions, decimals, and percents
\end{itemize}
\end{learningGoals}

\begin{conceptBox}{Circle Graphs}
A \textbf{circle graph} (pie chart) shows how a whole is divided into parts. The full circle represents $100\%$ (or $360°$).

\begin{itemize}[leftmargin=*]
    \item Each \textbf{sector} shows one category as a fraction of the whole.
    \item Sector angle $=$ percent $\times 360°$.
    \item All sectors must add up to $100\%$ (or $360°$).
\end{itemize}
\end{conceptBox}

\begin{workedExample}{Reading a Circle Graph}
A survey of $200$ students' favorite sport:

Soccer $40\%$, Basketball $25\%$, Baseball $20\%$, Other $15\%$

How many students chose soccer? $200 \times 0.40 = 80$

What angle for basketball? $0.25 \times 360° = 90°$
\answerHighlight{$80$ students chose soccer; basketball sector $= 90°$.}
\end{workedExample}

\mascotSays{To check your circle graph, make sure all the percents add to $100\%$!}

\begin{practiceBox}{Circle Graphs}
\resetProblems

\prob A circle graph shows that $30\%$ of $500$ people prefer cats. How many is that?
\answerExplain{$150$ people}{$500 \times 0.30 = 150$.}

\prob A sector represents $\frac{1}{4}$ of a data set. What is its angle?
\answerExplain{$90°$}{$\frac{1}{4} \times 360° = 90°$.}

\prob Four categories have percents $35\%$, $25\%$, $20\%$, and $x\%$. Find $x$.
\answerExplain{$20\%$}{Total must be $100\%$: $35 + 25 + 20 + x = 100$. So $x = 20\%$.}

\prob In a class of $40$ students, $12$ walk to school, $20$ take the bus, $8$ get a ride. What percent walks?
\answerExplain{$30\%$}{$\frac{12}{40} = 0.30 = 30\%$.}

\prob A sector has a $72°$ angle. What percent of the whole does it represent?
\answerExplain{$20\%$}{$\frac{72}{360} = 0.20 = 20\%$.}

\end{practiceBox}
